\documentclass[sigconf,review,anonymous]{acmart}
\AtBeginDocument{\providecommand\BibTeX{{Bib\TeX}}}
\setcopyright{none}
\settopmatter{printacmref=false}
\renewcommand\footnotetextcopyrightpermission[1]{}
\pagestyle{plain}
\usepackage{booktabs}
\usepackage{graphicx}
% Passages changed by the 2026-09-21 corrections are marked with a dagger; see the
% Corrections section at the end.
\newcommand{\corr}{\textsuperscript{\dag}}

\begin{document}

\title{Static Impact Analysis Does Not Transfer:\\A Pre-Registered Refutation of Two LLM-Agent Reliability Mechanisms}

\author{Anonymous Author(s)}

\begin{abstract}
Coding agents built on frozen language models are widely argued to need external
mechanisms that supply what the model structurally lacks: a memory of the codebase, a
prediction of what an edit will break, and a gate that refuses to act on an
under-specified request. We built two such mechanisms and reported strong results on
self-constructed evaluations: a dependency-graph \emph{impact oracle} with perfect recall
($1.000$) of impacted files, and a complexity \emph{router} with an assumption
\emph{gate} achieving $\mathrm{F1}=1.00$ and a $62.1\%$ cost saving. This paper reports
what happened when we evaluated both against data we did not construct, under a protocol
frozen before any data was collected.

Both results collapse. On $759$ files across nine independently selected open-source
Python repositories, with co-change ground truth mined from real commit history, the
oracle's pooled recall is $2.20\%$ (95\% CI $[1.20, 3.31]$) --- not $100\%$ --- and
$91.0\%$ of files receive \emph{exactly zero} recall. A grep baseline with no dependency
graph at all attains recall $0.573$ and F1 $0.437$ against the oracle's $0.042$. On $80$
held-out tasks drawn from real GitHub issues, with every threshold left untouched, gate
F1 falls from $1.00$ to $0.37$, and the cost saving does not merely shrink but inverts: counting the
pipeline's actual total spend, routing costs $20.2\%$ \emph{more} than always using the premium tier.

We report three things beyond the refutation. First, a structural decomposition of all
$18{,}221$ false negatives showing that $96.8\%$ have a static path in the oracle's own graph
that its reverse-only traversal cannot take --- a bound on what static analysis could reach, not
evidence that those misses are fixable\corr{} --- together with a construction defect whose safe
repair alone recovers an $11.0\times$ recall gap. Second, a repair of both defects evaluated
under the same frozen protocol with a declared tuning/held-out repository split. Third,
the observation that our central composition law is decades-old layer-of-protection
algebra, and was independently and more generally derived weeks before our own write-up.
We argue the transferable contribution is the protocol, not the mechanisms.
Corrections made on 2026-09-21 after an external review (listed at the end) tighten several
inferences --- the confidence intervals, the repaired oracle's comparison with grep, the labels,
and the ceiling --- and leave the refutation itself unchanged.
\end{abstract}

\begin{CCSXML}
<ccs2012>
<concept><concept_id>10011007.10011074.10011099</concept_id>
<concept_desc>Software and its engineering~Software verification and validation</concept_desc>
<concept_significance>500</concept_significance></concept>
<concept><concept_id>10011007.10011074.10011134</concept_id>
<concept_desc>Software and its engineering~Software evolution</concept_desc>
<concept_significance>300</concept_significance></concept>
</ccs2012>
\end{CCSXML}
\ccsdesc[500]{Software and its engineering~Software verification and validation}
\ccsdesc[300]{Software and its engineering~Software evolution}
\keywords{LLM coding agents, change impact analysis, co-change, negative results,
pre-registration, verification gates}

\maketitle

\section{Introduction}

A developer opening an unfamiliar file carries context the file does not contain. They
remember which callers depend on it, they anticipate what a signature change will break,
and when a request is too vague to act on they ask instead of guessing. A frozen language
model has none of this: it is a stateless map from a bounded context window to a token
sequence, and each invocation begins with no memory of the last.

This observation has motivated a large and growing class of \emph{external} mechanisms
that supply the missing faculties around an unmodified model --- persistent stores,
dependency-graph analyses, verification hooks, and routing layers. The argument for them
is structural and, we still believe, correct. What is far less clear is whether any
particular mechanism delivers what its authors measure.

We are those authors. We built two mechanisms and evaluated them the way such work is
commonly evaluated: on material we constructed ourselves. An \emph{impact oracle} parsed a
Python codebase into a dependency graph and answered ``what will this edit break'' by
reverse reachability with per-hop confidence decay; on a ten-file package with five
injected mutations it recovered every impacted file, recall $1.000$. A complexity
\emph{router} classified tasks into cost tiers and an assumption \emph{gate} halted
under-specified requests; on thirty hand-written tasks it achieved $\mathrm{F1}=1.00$ and
a $62.1\%$ cost saving against an always-premium baseline.

Both evaluations were honestly labelled as demonstrations. Neither survived contact with
data we did not build.

This paper reports that failure and what we learned from anatomising it. We pre-registered
a protocol --- corpus eligibility criteria, commit filters, thresholds, and bootstrap
seeds --- and froze it before cloning a single repository, with the explicit stipulation
that an unfavourable number would be reported rather than reconciled. We then evaluated
the oracle against co-change ground truth mined from nine independently selected
open-source projects, and the router and gate against tasks drawn from real GitHub issues,
with every threshold left exactly as shipped.

\paragraph{Contributions.}
\begin{enumerate}
\item A \textbf{pre-registered evaluation protocol} for LLM-agent impact mechanisms, using
repository co-change history as ground truth, with all thresholds and filters frozen before
data collection and an audited append-only amendment trail (\S\ref{sec:protocol}).
\item A \textbf{refutation} of both mechanisms' headline results on real data: oracle recall
$1.000 \rightarrow 0.022$, gate F1 $1.00 \rightarrow 0.37$, and a cost saving that
inverts to $-20.2\%$ once escalation retries are counted (\S\ref{sec:eval1}, \S\ref{sec:eval2}).
\item A \textbf{structural failure-mode taxonomy} showing that $96.8\%$ of missed impact has a
static path the traversal cannot take, which bounds what a static method could reach\corr{}, and isolating a
graph-construction defect whose safe repair recovers an $11.0\times$ recall gap
(\S\ref{sec:failure}, \S\ref{sec:repair}).
\item A \textbf{repair and re-evaluation} of both defects under the same frozen protocol,
with a declared tuning/held-out repository split (\S\ref{sec:repair}).
\item An \textbf{honest novelty accounting}: our composition law is an instance of
layer-of-protection analysis, and was independently and more generally derived weeks
before our own write-up (\S\ref{sec:related}).
\end{enumerate}

\paragraph{What we do not claim.} We do not claim that external impact analysis is
worthless for coding agents, that co-change is a substitute for semantic impact, or that
our repaired mechanism is production-ready. We claim that the evidence standard commonly
used to justify such mechanisms --- a self-built demonstration --- is capable of
overstating recall by more than an order of magnitude, and we demonstrate this on our own
work.


\section{Protocol}\label{sec:protocol}

The evaluation was pre-registered. Before any repository was cloned or any commit mined, we
fixed and froze: nine numbered corpus eligibility criteria; the commit filters defining
ground truth; the confidence-threshold grid; the bootstrap procedure and its seed; and the
stipulation that governs everything else --- \emph{if a downstream number looks bad, the fix
is to report it, not to edit the protocol}. Amendments were permitted only as dated,
append-only addenda. Three were filed, all documenting the corpus-selection funnel; a
character-level diff confirms nothing above the freeze line was altered.

\paragraph{Corpus.} Candidates required Python purity $\geq 0.85$, $\geq 3000$ stars,
$\geq 50$ contributors, $\geq 800$ commits, age $\geq 5$ years, activity within 24 months,
repository size under 500\,MB, no affiliation with this project, and $\geq 15$ Python files
after path exclusions. The grep baseline extends the prototype's own symbol-search routine to file
granularity; matching is on identifier occurrences, so short or generic identifiers can inflate its
false-positive count, which if anything favours the oracle in the comparison. Twelve primary candidates were tried in a pre-declared order with
substitution from a pre-declared backup list; substitutions were permitted only in response
to these metadata criteria, never in response to performance. Nine repositories qualified:
\texttt{requests}, \texttt{flask}, \texttt{click}, \texttt{jinja}, \texttt{rich},
\texttt{httpx}, \texttt{attrs}, \texttt{pytest}, and \texttt{loguru}. Each was pinned to an
exact commit SHA recorded in the corpus artifact.

\paragraph{Ground truth.} For a file $f$, the ground-truth impact set is the set of other
Python files appearing in the same non-merge commit as $f$, aggregated over history. Merges,
mega-commits above the pre-declared size threshold (formatting sweeps, vendored dependency
bumps, mass renames), and excluded paths were dropped by the frozen filters. This yields
801 labelled files, of which 759 were evaluated: the pre-registered cap of 200 files per
repository cut \texttt{pytest} from 242 to 200\corr{}.

\paragraph{Ground truth is a proxy, and its ceiling is measurable.} Co-change is not
semantic impact. Files co-change for reasons no static analysis can see (release chores,
coordinated documentation, test-fixture coupling), and files that \emph{should} co-change
sometimes do not. We therefore measured the ceiling directly: the fraction of ground-truth
pairs connected by \emph{any} static path in either direction, at any hop, in the graph the
oracle itself builds. Pooled, 96.9\% of ground-truth impact is statically reachable in
principle. That number bounds every result below, and it is what makes the observed recall interpretable
rather than merely low. Two caveats: it is measured on the graph the \emph{as-shipped} oracle builds, so
repairing the construction defect (\S\ref{sec:repair}) raises it slightly, to roughly $97.2\%$; and
reachability in a dense graph is a weak property, so the ceiling is an upper bound on what any static
method could attain, not evidence that a reachable pair is causally related.

\paragraph{Independent reproduction.} Before computing any metric, we re-mined all nine
repositories' histories from scratch and compared against the supplied ground truth: 801 of
801 rows matched exactly. Two harness corrections were made at this stage --- both
identified and disclosed \emph{before} any metric was computed.

\section{Evaluation I: Impact Closure}\label{sec:eval1}

\begin{table}[t]
\caption{Pooled, micro-averaged over 759 files in 9 repositories, with 95\% bootstrap CIs
($B=2000$, file-resampled). The oracle is shown at $t=0.02$, the most favourable point of
the pre-registered grid for its recall. File-resampled intervals are too narrow for this
design; repository-cluster intervals are given in the Corrections\protect\corr{}.}
\label{tab:eval1}
\small
\begin{tabular}{lccc}
\toprule
Method & Precision & Recall & F1 \\
\midrule
Impact oracle ($t{=}0.02$) & 0.398 & \textbf{0.022} & 0.042 \\
Impact oracle (shipped $t{=}0.10$) & 0.381 & 0.017 & 0.033 \\
Grep baseline (no graph) & 0.353 & \textbf{0.573} & \textbf{0.437} \\
Edited-file-only & 0.000 & 0.000 & 0.000 \\
\midrule
\emph{Static ceiling on this ground truth} & --- & \emph{0.969} & --- \\
\bottomrule
\end{tabular}
\end{table}

Table~\ref{tab:eval1} states the result. The oracle's pooled recall is 0.0220, 95\% CI
[0.012, 0.033] when files are resampled, and [0.0005, 0.052] when the nine repositories are
resampled instead\corr{}. The claim under test was $1.000$. The oracle reaches 2.27\% of its own
static ceiling.

The per-file distribution shows this is not a uniform shortfall but close to a step
function: 691 of 759 files (91.0\%) receive \emph{exactly zero} recall --- the
oracle misses every co-changed file --- and one file achieves recall $1.0$. A mean of $2.2\%$
conceals a system that, for nine files in ten, returns nothing useful at all.

The grep baseline is the uncomfortable comparison. A textual search for symbol names, with
no dependency graph, no confidence model, and no persistent state, attains recall 0.573
and F1 0.437 against the oracle's 0.042. Whatever a codebase world-model is supposed to
buy over an unstructured baseline, it is not visible here. Precision is the one metric where
the oracle is competitive (0.398 versus 0.353): when it warns, it is often right. It
simply almost never warns.

\begin{figure*}[t]
\centering
\includegraphics[width=\textwidth]{figures/fig_cochange_results.png}
\caption{Per-repository recall, precision, and F1 with 95\% bootstrap CIs. The oracle's
recall bars are near-invisible against the grep baseline in every repository.}
\label{fig:eval1}
\end{figure*}

\section{Failure Modes}\label{sec:failure}

A negative result is only useful if it explains itself. We classified all 18221 pooled false
negatives structurally, by testing reachability in the oracle's own graph.

\begin{table}[t]
\caption{Structural decomposition of all 18221 false negatives.}
\label{tab:failure}
\small
\begin{tabular}{lcp{3.1cm}}
\toprule
Category & Share & Interpretation \\
\midrule
Sibling (common dependency) & 94.68\% & Static path exists\corr{}: both files depend on a shared module;
reachable only by a mixed-direction path \\
Forward-only dependency & 2.13\% & Static path exists\corr{}: excluded by design \\
No static path at all & 3.19\% & True ceiling: ground-truth limit \\
\bottomrule
\end{tabular}
\end{table}

The decomposition in Table~\ref{tab:failure} is the paper's most actionable finding.
Only 3.19\% of missed impact is attributable to the proxy nature of co-change. The
remaining 96.8\% is a path the oracle's own graph already contains but its
traversal cannot reach: 94.7\% are \emph{sibling} pairs, where $A$ and $B$ both import a
common module $C$ and co-change when $C$'s contract shifts, yet neither is reverse-reachable
from the other. The mechanism was not defeated by the messiness of real code; it was defeated
by answering a narrower question than the one that matters.

Two cautions apply to this decomposition\corr{}. Co-change is symmetric while impact is
directional, so a pair with a static path in either direction is not thereby a pair whose impact a
static method should predict. And we ran no control measuring how many \emph{random} file pairs
are connected in the same dense graph, so the $96.8\%$ bounds what static analysis could reach; it
does not show that those misses are fixable.

\paragraph{A construction defect.} Separately, and reported as a labelled diagnostic rather
than substituted into the numbers above, we found a bug. The world-model's module-qualified-name
derivation collides with absolute imports in \texttt{src/}-layout packages, creating
disconnected phantom nodes. A guard exists for precisely this case but never fires: its
membership test is defeated by the graph library's edge-insertion auto-vivification, which
creates the node before the guard inspects it. This explains why recall was approximately zero
in exactly the six \texttt{src}-layout repositories and non-negligible only in the three
flat-layout ones. A naive repair of the guard raises pooled recall by $14.5\times$, but is unsafe
(\S\ref{sec:repair}); the safe repair we actually apply recovers $11.0\times$, and that is the
figure we report throughout.

\section{Evaluation II: Routing and Gating}\label{sec:eval2}

We scored the router and gate on 80 tasks drawn from real GitHub issues and pull requests,
with independent gold labels for cost tier and for whether the request is under-specified
enough that an agent should halt. \textbf{No threshold, weight, or pattern was modified.}

\paragraph{Label noise first.} Agreement between two labelling passes was
$\kappa = 0.952$ (quadratic-weighted) for tier and $\kappa = 0.516$ for the halt
decision. Tier labels are near-perfect; halt labels are moderate. Both passes were produced by
the same model (\texttt{claude-sonnet-5}) with differently worded prompts, on a 30-task subset,
and that model is also the correctness judge and the mid-tier executor, so $\kappa$ measures how
robust the labels are to prompt wording, not whether they are valid\corr{}. We report this before the
accuracy numbers because it bounds how much of the degradation below can be attributed to
label ambiguity rather than to the mechanism --- and for the halt decision, some of it can.

\begin{table}[t]
\caption{Tuned 30-task set versus 80 held-out real-issue tasks, thresholds unchanged.}
\label{tab:eval2}
\small
\begin{tabular}{lcc}
\toprule
Metric & Tuned ($n{=}30$) & Held-out ($n{=}80$) \\
\midrule
Gate F1 & 1.00 & \textbf{0.37} \\
Gate precision & 1.00 & 0.45 \\
Gate recall & 1.00 & \textbf{0.31} \\
Routing exact-tier & 1.00 & 0.53 \\
Routing within-1-tier & 1.00 & 0.91 \\
Cost saving & 62.1\% & \emph{see Table~\ref{tab:cost}} \\
\bottomrule
\end{tabular}
\end{table}

Gate F1 falls from $1.00$ to 0.37, driven by recall collapsing to 0.31: the gate fails
to flag 11 of the 16 genuinely under-specified issues. Exact-tier routing accuracy falls to
0.53, though within-one-tier accuracy is comparatively robust at 0.91 --- most routing
errors are off-by-one rather than wild. Cost requires two orthogonal distinctions the tuned demonstration never had to make, because on that set
almost every routed output was correct. \emph{Framing}: whether we count only the first, initially-routed
attempt, or everything the pipeline actually spent including each escalation retry. \emph{Gating}: whether
a dollar saved counts when the cheaper output is wrong. Table~\ref{tab:cost} gives all four cells.

\begin{table}[t]
\caption{Held-out cost saving versus an always-premium baseline, by framing and gating (95\% bootstrap CIs).
The tuned demonstration's $62.1\%$ corresponds to the top-left cell: first-attempt, ungated.}
\label{tab:cost}
\small
\begin{tabular}{lcc}
\toprule
& Ungated & Correctness-gated \\
\midrule
First attempt only & $59.5\%$ & $3.6\%$ \\
& \tiny[50.5, 68.2] & \tiny[0.0, 8.5] \\
Total pipeline spend & $\mathbf{-20.2\%}$ & $2.8\%$ \\
& \tiny[$-27.5$, $-12.3$] & \tiny[$-0.8$, 7.8] \\
\bottomrule
\end{tabular}
\end{table}

Read in the order that matters. Routing does identify that most of these tasks do not need the premium
tier, so on first attempts alone it saves $59.5\%$ in raw dollars --- close to the demonstration's
$62.1\%$, and the one cell where the tuned result survives. But almost none of that cheaper output is
\emph{correct}: gating the same first-attempt framing on the judge's verdict collapses the saving to
$3.6\%$. And the pipeline does not stop at the first attempt. Counting what it actually spent chasing
correctness up the ladder, the raw saving is $\mathbf{-20.2\%}$: the router costs \emph{more} than always
sending every task to the premium tier, because it pays for cheap, then mid, then premium in sequence and
still mostly fails. The gated total-spend figure, $2.8\%$, has a confidence interval spanning zero.

The honest summary is not that the saving shrank but that it inverted. On this task distribution a
complexity router is a net cost \emph{increase} once its retries are counted, and the single headline
number the demonstration produced ($62.1\%$) turns out to be the most favourable of four available cells.

The inversion is also largely mechanical\corr{}. On 58 of the 64 tasks no tier produced an output
the judge accepted, and always-premium was accepted on only 3 of 64, so on almost every task the
pipeline paid for every tier it tried. Measured per accepted output, the pipeline cost \$1.06
against always-premium's \$1.76, about 40\% less --- but from 6 and 3 accepted outputs, so neither
figure is stable, and we report it beside the $-20.2\%$ rather than instead of it.

\paragraph{The gate's confidence score is non-monotonic.} Binned calibration is not monotone in
the score: decision accuracy is $0.93$ in the middle bin ($n{=}28$) and $0.94$ in the highest
($n{=}16$), but dips to $0.50$ in the second-highest bin ($n{=}4$). The bins were meant to be
equal-frequency, but scores tie at 0.63, 0.90 and 1.00, so they hold 27, 5, 28, 4 and 16 tasks;
the expected calibration error is 0.103 with half-open $(lo, hi]$ bins and 0.078 with $[lo, hi)$
bins. We earlier called the dip ``well within sampling noise''. It is not, at the nominal level: if
the bin's accuracy matched the middle bin's $0.93$, two or fewer correct of four has probability
$0.028$\corr{}. With four tasks and no correction for the bins we looked at, it is weak evidence
either way. What it does establish is that the score is not
reliably ordered across its range, so choosing an operating point from it requires more evidence
than this sample provides.

\paragraph{Verification coverage.} Of the 64 well-specified held-out tasks, \emph{none} admitted
execution-based verification: real issue text does not come with runnable input/output examples.
Correctness was therefore judged by a model-based criterion, which is weaker, and we mark the
cost figure accordingly. One candidate execution match was found by pattern search and rejected
on inspection as prose describing a source-code branch rather than an invocable example.


\section{Repair and Re-Evaluation}\label{sec:repair}

A refutation that stops at the negative number leaves the more useful question unasked: was the
mechanism wrong, or merely built wrong? \S\ref{sec:failure} says $96.8\%$ of the missed impact is
reachable in a graph the tool already constructs. We therefore repaired both defects and re-ran the
same frozen protocol.

\paragraph{Honest parameter setting.} The traversal repair introduces new parameters, so it cannot be
evaluated on the corpus that motivated it --- that is the error this paper is about. Before writing
any repair code we declared a stratified split: six \textbf{tuning} repositories (\texttt{rich},
\texttt{httpx}, \texttt{click}, \texttt{flask}, \texttt{jinja}, \texttt{attrs}) and three
\textbf{held-out} repositories (\texttt{loguru}, \texttt{requests}, \texttt{pytest}), stratified by
package layout in proportion to the corpus. All eight parameters were fixed by grid search on the
tuning repositories only, written to a frozen parameter file, and never adjusted afterwards. The
held-out repositories were touched exactly once, after freezing.

\paragraph{Defect 1: construction.} We repaired the phantom-node merge --- but not as the diagnostic
suggested. The naive guard fix is \emph{unsafe}: it merges a bare \texttt{import json} into an
unrelated local \texttt{pkg.json} submodule, fabricating dependency edges that do not exist. We ground
matching in the layout prefix instead. This yields a smaller recall improvement than the unsafe
variant ($11.0\times$), and we prefer it: a tool that invents edges to raise recall is worse than one
that misses them.

\paragraph{Defect 2: traversal.} We added terminal sibling and forward relations with per-relation
weights and hop caps. The sibling relation must be capped by the shared dependency's in-degree ---
admitting every sibling of a high-traffic utility module implicates most of the codebase. A
post-freeze sensitivity check on the tuning repositories confirms the cap does real protective work
rather than passing through unbounded.

\begin{table}[t]
\caption{Repair sequence under the frozen protocol. Row (d) is the headline: three repositories never
seen during parameter selection. Rows (a)--(b) are pooled over all nine.}
\label{tab:repair}
\small
\begin{tabular}{llccc}
\toprule
& Configuration & P & R & F1 \\
\midrule
(a) & As shipped (reproduced exactly) & 0.398 & 0.022 & 0.042 \\
(b) & \quad + construction fix & 0.304 & 0.242 & 0.270 \\
(c) & \quad + traversal fix, \emph{tuning} & 0.461 & 0.668 & 0.545 \\
(d) & \quad + traversal fix, \textbf{held-out} & 0.305 & 0.653 & \textbf{0.416} \\
(d$'$) & \quad\quad same, at tuned $t{=}0.10$ & 0.320 & 0.647 & 0.428 \\
\midrule
& grep baseline, held-out & 0.269 & 0.601 & 0.371 \\
& \emph{static ceiling, held-out repos} & --- & \emph{0.977} & --- \\
\bottomrule
\end{tabular}
\end{table}

\paragraph{Result.} We headline row (d) at the canonical $t{=}0.02$ used for the as-shipped baseline, so the
comparison is like-for-like. On repositories never seen during parameter selection the repaired oracle reaches
$\mathrm{F1}=0.416$ (95\% CI $[0.382, 0.446]$) against grep's 0.371 (CI $[0.332, 0.405]$). The point
estimate reverses the as-shipped ordering, in which grep won by roughly an order of magnitude; whether
the repaired oracle actually beats grep is not established (below)\corr{}. Recall rises from 0.022 to 0.653,
reaching $66.8\%$ of the held-out repositories' own static ceiling ($0.977$), and the share of files with
exactly zero recall falls from $91.0\%$ to $19.3\%$. Row (d$'$) reports the same configuration at
$t{=}0.10$, the F1-peak of the tuning grid; it scores marginally higher, and we decline to headline it
because that threshold was selected on the tuning repositories and is therefore a ninth tuned parameter
rather than a pre-registered one. All 49/49 tests pass, including 13 new tests covering
both fixes and the stdlib-collision safety case. Both fixes are necessary: neither alone crosses the
baseline.

\paragraph{The advantage is not established, and precision generalised worse than recall.} We
earlier read the comparison as a real but modest win\corr{}. The evidence does not support that.
The two file-level intervals overlap. The paired difference in F1 is $+0.044$ at $t{=}0.02$, and
$+0.057$ at $t{=}0.10$, where per-repository counts exist; there all three held-out repositories
favour the oracle ($+0.114$ \texttt{loguru}, $+0.131$ \texttt{requests}, $+0.033$ \texttt{pytest}),
but three out of three gives a one-sided sign-test $p$ of $0.125$, and \texttt{pytest} supplies
71\% of the held-out pairs. The split was also clean only for the eight numeric parameters: the
choice of which relations to add came from the failure analysis of \S\ref{sec:failure}, pooled
over all nine repositories, held-out ones included. More instructive is the asymmetry: recall transferred almost perfectly from tuning to held-out repositories
($-0.021$), while precision fell by $-0.141$. Some overfitting to the tuning set occurred during
parameter selection --- visible \emph{only} because the split was declared in advance. Had we tuned on
all nine repositories we would have reported the better number and never known the gap existed.

\begin{figure*}[t]
\centering
\includegraphics[width=\textwidth]{figures/fig_repair_beforeafter.png}
\caption{The repair sequence (a)$\rightarrow$(d) with 95\% bootstrap CIs, against the grep baseline
and the static ceiling.}
\label{fig:repair}
\end{figure*}


\section{Threats to Validity}\label{sec:threats}

\paragraph{Construct.} Co-change is a proxy for semantic impact and errs in both directions.
Files co-change for reasons no dependency analysis can predict --- release chores, coordinated
documentation, test-fixture updates --- and files that ought to co-change are sometimes missed
by their authors. We do not treat co-change as truth; we measure its ceiling (96.9\% of
ground-truth pairs are statically connected in principle) and interpret every result against
it. An over-warning is likewise not necessarily an error: it may be a correct dependency that
has not yet co-changed in the mined window.

\paragraph{Construct, continued: the correctness judge.} No held-out task admitted execution-based
verification (0/64; the mechanical check was re-run independently and its single candidate hit rejected on
inspection), so correctness was scored by a fixed-rubric model judge. The judge was pre-declared as
``neither cheap nor premium'' --- neutral with respect to the \emph{tier ladder}, but not neutral in an
absolute sense, because the ladder's mid tier is that same model. For the 38 of 64 tasks initially routed
to mid, and for every mid-tier escalation attempt, the judge grades output from its own model family.
Evaluators are known to favour their own generations, which biases this measurement in the direction of
\emph{overstating} mid-tier correctness and therefore overstating the routed pipeline's savings; the
inversion we report is if anything conservative. Separately, seven generation calls returned empty output
at the pre-declared token budget and were retried at the same tier with a raised budget; this is logged as
a deviation rather than silently patched, and no escalation cost is attributed to those retries.

\paragraph{Internal.} The pre-registration, the tuning/held-out split, and the parameter freeze were all
self-administered within one continuous working session, with no external timestamping authority: a reader
must take the ordering on trust, evidenced only by the artifact store's own version timestamps and the
append-only amendment trail. We consider this the central weakness of a self-evaluation and state it plainly
rather than resting on the protocol's authority. The held-out task labels come from a single model,
queried twice with differently worded prompts\corr{}; halt-decision
agreement is only moderate ($\kappa = 0.52$), so part of the gate's measured degradation is
label ambiguity rather than mechanism failure. Tier agreement is near-perfect ($\kappa = 0.95$),
so the routing result is on firmer ground. No held-out task admitted execution-based
verification, so correctness was judged by a weaker model-based criterion.

\paragraph{External.} Nine repositories, Python only, open-source only, all library-shaped
rather than application-shaped. Results may differ for applications with heavy dynamic
configuration, for languages with static typing that makes dependencies more explicit, or for
proprietary codebases with different commit conventions. GitHub issue text is a proxy for the
prompts a developer actually gives an agent.

\paragraph{Conclusion.} All intervals are 95\% file- or task-resampled bootstrap intervals.
Files are not independent units here --- every ground-truth pair is counted from both of its
files, and files cluster in nine repositories --- so the file-resampled intervals are too narrow;
repository-cluster intervals are in the Corrections\corr{}.
We report them for every headline number and make no claim that survives only at a point
estimate. We apply no multiple-comparison correction and claim no significance test; the
comparisons are descriptive.

\section{Related Work}\label{sec:related}

\paragraph{The composition law is not ours.} Read as mathematics, the product-of-independent-layer-misses
form is textbook. Layer-of-protection analysis has expressed unmitigated event frequency as a
product over independent protection layers since the 1990s~\cite{dowell1998}; fault-tree analysis
generalises it~\cite{ruijters2015}; Reason's Swiss-cheese model gives the same intuition a
human-factors vocabulary~\cite{reason1990}; and reliable-computation-from-unreliable-parts runs
back to von Neumann and to N-version programming~\cite{avizienis1985}. The literature also supplies
the caveat any layered argument must confront: Knight and Leveson showed experimentally that
independently written versions do \emph{not} fail independently~\cite{knight1986}, so an
independence assumption between layers must be defended rather than asserted. Our own formal
synthesis asserted it: it multiplied the catch probabilities of one classifier run at three
lifecycle points on the same diff, which understates the residual 400-fold in its own worked
example\corr{}. We claim no novelty for the arithmetic.

\paragraph{And it was independently derived, more generally, before ours.} More consequentially,
two concurrent preprints derive a strictly more general Bayesian form of the same law weeks before
our own write-up. \emph{Odds Law}~\cite{aksu2026} develops a decomposition algebra whose
verification-odds law and $k$-gate geometric-amplification theorem subsume our two-layer product as
a special case, adding a threshold dichotomy and an information-theoretic ceiling we do not have.
Its companion \emph{Maestro Order}~\cite{aksu2026a} compiles the algebra into a controller and names
the coder/test-suite pairing as its canonical software-engineering instance --- the same
probabilistic-generator/deterministic-checker identification we had claimed. We concede priority.
What our framing retains is narrower: an argument for the split \emph{derived} from the architecture
of a frozen parametric map rather than assumed as a black-box parameter, and --- the gap we can
actually fill --- an evaluation against real repositories, where both preprints offer only Monte
Carlo simulation of an abstract solver model.

\paragraph{Impact analysis is forty years old and being reinvented right now.}
The oracle's pipeline descends directly from program slicing~\cite{weiser1984}, program dependence
graphs~\cite{ferrante1987}, and interprocedural slicing~\cite{horwitz1990}. Its ground truth follows
the co-change mining line~\cite{zimmermann2005,zimmermann2004}, and its confidence decay has a
decade-old empirical precedent in dependency-distance studies of change propagation. Concurrently,
at least four 2026 preprints wire comparable AST/graph analyses into LLM coding-agent loops ---
source-to-test dependency mapping, library-defect propagation, persistent codebase knowledge graphs,
and assumption localisation --- and, notably, \emph{none of them claims perfect recall}. That the
sub-field is simultaneously decades old and actively re-invented raises the evidential bar rather
than lowering it.

\paragraph{Pre-registration in software engineering.} Pre-registration and Registered Reports are not new
to empirical software engineering, where they were introduced to address exactly the degrees of freedom we
exploit here unintentionally. We claim no novelty for the practice. What we add is narrower: a concrete
instance of applying it to LLM-agent mechanism evaluation, where self-built demonstrations are currently the
prevailing evidence standard, together with a measurement of how large the resulting gap can be.

\paragraph{Routing and abstention.} Cost-aware cascading was named by
FrugalGPT~\cite{chen2023} and the successor literature is large, with learned and preference-trained
routers now standard. The most pointed prior work for us argues that deployed routers rely on
uncalibrated confidence and per-workload threshold tuning, and replaces hand-set thresholds with
calibrated error probabilities --- a critique that lands squarely on our own tuned-rubric
methodology, and which our held-out result independently confirms. Selective prediction with a
reject option is older still, running back to Chow's rule. Surveys of LLM self-correction report
increasingly negative results, which supports rather than undermines the case for an
\emph{external} deterministic check.

\section{Discussion}\label{sec:discussion}

Three lessons generalise beyond our two mechanisms.

\textbf{A self-built demonstration can overstate recall by more than an order of magnitude.}
Our demo package was not adversarial or dishonest: ten files, five plausible mutations, honestly
labelled as a demonstration. It still reported $1.000$ where real repositories gave $0.022$. The
demo was unrepresentative in a specific and instructive way --- it used a flat package layout, and
the defect that destroyed recall only manifests under \texttt{src/} layouts. A demonstration
selects its own structure, and structure was the variable that mattered.

\textbf{Tuning and evaluating on the same set does not degrade gracefully.} The gate did not fall
from $1.00$ to $0.85$; it fell to $0.37$, and its confidence score is non-monotonic in the region
where an operating point would be chosen. The failure is qualitative, not a matter of degree.

\textbf{Report the ceiling, not just the score.} Decomposing false negatives structurally was what
turned an uninterpretable $2.2\%$ into an actionable finding: $96.8\%$ with a static path the
traversal could not take, $3.2\%$ with none. The first figure bounds what a repair could reach; it
does not show those misses are fixable\corr{}. Without the ceiling, the number is merely discouraging; with it, it is a specification
for the repair.

\section{Conclusion}\label{sec:conclusion}

We pre-registered an evaluation of two LLM-agent reliability mechanisms we had built ourselves and
reported strong results for, then ran it against data we did not construct. Both headline results
collapsed: recall $1.000 \rightarrow 0.022$, gate F1 $1.00 \rightarrow 0.37$, and a cost saving
that inverts from $+62.1\%$ to $-20.2\%$ once the pipeline's escalation retries are counted. A grep baseline with no dependency graph
outperformed the graph-based oracle. We diagnosed the causes structurally, repaired both defects,
and re-evaluated under the same frozen protocol with a declared held-out split.

We think the protocol is the transferable contribution. The mechanisms may yet be made to work ---
$96.8\%$ of the missed impact has a static path in a graph the tool already builds\corr{} --- but the evidence
standard that told us they already worked was the thing most in need of repair.

\section*{Corrections (2026-09-21)}

An external deep review of the repository that hosts this paper (2026-09-21) recomputed its
statistics from the replication package. Every reported count reproduced exactly. Some inferences
drawn from them were too strong. The passages marked \dag{} above were changed; this list gives
each change. All figures were recomputed with \texttt{research/recompute\_corrections.py}
(standard-library Python, fixed seeds).

\begin{enumerate}
\item \textbf{Confidence intervals.} Every ground-truth pair is mirrored --- all 20,144 pairs among
the 801 labelled files have both ends labelled and are counted from both sides --- and files
cluster in nine repositories, so file-level resampling understates uncertainty. Resampling
repositories (seed 1234, 20,000 resamples), the oracle's pooled precision interval is
$[0.15, 0.91]$ (file-level $[0.27, 0.57]$) and its recall interval $[0.0005, 0.052]$ (file-level
$[0.012, 0.033]$). The refutation survives: oracle F1 $[0.001, 0.093]$ against grep's
$[0.381, 0.539]$, and the paired difference (grep minus oracle) is $[0.342, 0.517]$. With nine
clusters the interval ends move in the third decimal from seed to seed.
\item \textbf{801 and 759.} 801 files were labelled; 759 were evaluated, because the pre-registered
cap of 200 files per repository cut \texttt{pytest} from 242 to 200. The text had said 759 ``met
the evaluation's inclusion condition''.
\item \textbf{The repaired oracle does not demonstrably beat grep.} The original text read the
held-out comparison as ``a real but modest advantage''. See \S\ref{sec:repair} for the paired
differences, the sign test ($p = 0.125$), the 71\% \texttt{pytest} share, and the structural
choice that was made on all nine repositories.
\item \textbf{Labels.} The gold labels and the ``independent'' second pass were the same model with
two prompts, on $n = 30$; the same model is the judge and the mid-tier executor. $\kappa$ measures
prompt robustness, not label validity.
\item \textbf{Cost per accepted output.} The paper reported the $-20.2\%$ inversion without the cost
per output the judge accepted: \$1.06 for the pipeline against \$1.76 for always-premium, from 6
and 3 accepted outputs of 64. 58 of 64 tasks failed at every tier.
\item \textbf{The ceiling.} ``$96.8\%$ fixable'' treated symmetric co-change as directional impact
and had no random-pair control. It now reads as a bound on what static analysis could reach.
\item \textbf{Calibration.} The bins held 27, 5, 28, 4 and 16 tasks, not equal numbers; ECE is
0.103 or 0.078 depending on how ties at the bin edges are assigned; and the dip in the
second-highest bin is not ``well within sampling noise'' ($p = 0.028$ against the middle bin's
accuracy).
\item \textbf{The composition law.} Related work now notes that our own synthesis asserted the
independence between layers that Knight and Leveson warn against.
\end{enumerate}

The compiled PDF and the copy of this source inside \texttt{replication\_package.tar.gz} predate
these corrections.

\section*{Artifact availability}

All materials are available: the pre-registration with its append-only amendment trail, the corpus manifest
with pinned commit SHAs, the mined co-change ground truth, the held-out task set with its labelling protocol
and both labelling passes, the tuning/held-out split declaration and frozen parameter file, the repaired
package with its test suite, and the analysis code that regenerates every number and figure in this paper.
References are withheld for anonymous review and will be released on acceptance.

\section*{Design provenance}

The design principles behind these mechanisms --- verify a report before acting on it, do not act
without knowledge of consequences, treat a delegated capability as a trust --- were derived from an
ethical-epistemic reading documented in full in the extended version of this work. That reading
motivated which safeguards were considered mandatory rather than optional; it makes no technical
claim, and every result above stands or falls on the engineering evidence alone.

\bibliographystyle{ACM-Reference-Format}
\bibliography{refs}

\end{document}
